\section*{Aktivitas Pembelajaran dan Praktikum}
\addcontentsline{toc}{section}{Aktivitas Pembelajaran dan Praktikum}

\begin{table}[H]
\centering
\caption{Rencana aktivitas Bab X}
\label{tab:ch10-aktivitas}
\begin{tabularx}{\textwidth}{@{}llX@{}}
\toprule
\textbf{Sesi} & \textbf{Bentuk} & \textbf{Kegiatan} \\
\midrule
1 & Ceramah & Recursive descent dan predictive parsing \\
2 & Latihan & Menghitung FIRST dan FOLLOW \\
3 & Lokakarya & Membangun tabel LL(1) dan menelusuri stack \\
4 & Praktikum & M2: parser dan pengujian \\
\bottomrule
\end{tabularx}
\end{table}

\subsection*{M2 / Praktikum 2: Parser}
\textbf{Tujuan.} Membangun recursive descent parser untuk Subset-C
yang menghasilkan AST dan menangani kesalahan sintaksis.

\textbf{Keluaran.}
\begin{enumerate}[label=\arabic*.,leftmargin=*]
  \item Parser untuk ekspresi, deklarasi, dan pernyataan kontrol.
  \item AST minimal untuk ekspresi dan pernyataan.
  \item Minimal 20 kasus uji: valid dan tidak valid.
  \item Diagnostik sintaksis dengan lokasi dan harapan token.
  \item Laporan singkat.
\end{enumerate}

\textbf{Langkah kerja.}
\begin{enumerate}[label=\arabic*.,leftmargin=*]
  \item Hitung FIRST dan FOLLOW untuk grammar hasil Bab IX.
  \item Bangun tabel LL(1) dan periksa konflik.
  \item Implementasikan fungsi parser per nonterminal.
  \item Bangun AST saat mengurangi aturan.
  \item Tambahkan sinkronisasi dan pesan kesalahan.
  \item Uji dengan program valid dan tidak valid.
\end{enumerate}

\textbf{Rubrik.}
\begin{table}[H]
\centering
\caption{Rubrik Praktikum 2}
\label{tab:ch10-rubrik}
\begin{tabularx}{\textwidth}{@{}lXr@{}}
\toprule
\textbf{Aspek} & \textbf{Indikator} & \textbf{Bobot} \\
\midrule
Kebenaran & Parser menerima program valid, menolak yang tidak valid &
40\% \\
Kualitas kode & Struktur fungsi, penamaan, komentar & 30\% \\
Kelengkapan & Ekspresi, pernyataan, AST, diagnostik & 20\% \\
Pengujian & Kasus normal, batas, dan kesalahan & 10\% \\
\bottomrule
\end{tabularx}
\end{table}

\textbf{Definisi selesai.}
\begin{checklist}
  \item Parser lulus 20 kasus uji.
  \item AST terbentuk untuk program valid.
  \item Diagnostik menyebut lokasi dan harapan.
  \item Laporan diserahkan.
\end{checklist}
